% The mechanism, as one figure.
%
% Hand-written rather than generated, because it carries no numbers: it is the geometry a
% reader has to hold in mind, and geometry is exactly what prose conveys badly. The
% distinction it exists to make is that the slots are rank blocks *inside* one adapter, not
% one adapter each -- the neighbouring work that organises adapters into slots does the
% latter, and a reader who skims will assume we do too.
%
% Grayscale-safe by construction: every distinction is fill, dash or line weight, never
% colour. No font is set on the picture itself, so labels inherit the document font.
%
% Layout rules, each learned from a rendered version that broke them:
%   * every band of content owns a strip of y-space with at least 0.4cm of clearance from
%     the band above, because a label anchored "above" a shape grows upward into it;
%   * boxes are wide enough for their text at \footnotesize, since a clipped line inside a
%     box looks like a rendering fault rather than a mistake in the figure;
%   * the "$\cdot$" between $A$ and $B$ goes in the gap between them, never at a coordinate
%     that happens to fall inside a block.
\begin{figure}[t]
\centering
\begin{tikzpicture}[x=1cm,y=1cm,every node/.style={font=\small}]

  % ---- band 1 (y = 5.58 to 6.53): where the side path hangs off ----------------------
  \node[draw,line width=0.7pt,minimum width=3.6cm,minimum height=0.95cm,fill=black!7]
    (backbone) at (2.40,6.05) {frozen backbone $M_\theta$};
  \node[draw,line width=0.7pt,minimum width=4.6cm,minimum height=0.95cm,fill=white]
    (layer) at (9.20,6.05) {one projection layer $W_0$};
  \draw[-{Latex[length=2.2mm]},line width=0.8pt] (4.35,6.05) -- (6.75,6.05);
  \node[above=0.05cm,font=\footnotesize] at (5.55,6.05) {unchanged};

  % ---- band 2 (y = 3.85 to 5.15): the adapter and its rank split ---------------------
  \node[anchor=east,font=\small] at (3.05,4.45) {$W = W_0 + \frac{\alpha}{r}\,B\,A$};
  \draw[-{Latex[length=2.2mm]},line width=0.7pt] (3.15,4.45) -- (3.95,4.45);

  % A: one strip of K contiguous rank blocks; block 2 is the slot in play (filled).
  \foreach \i/\shade in {1/white,2/black!55,3/white,4/white} {
    \draw[line width=0.8pt,fill=\shade] ({3.95+1.30*(\i-1)},3.90) rectangle ++(1.30,0.72);
  }
  \foreach \i in {1,2,3,4} {
    \node[below=0.04cm,font=\footnotesize] at ({4.60+1.30*(\i-1)},3.90) {$A_{\i}$};
  }
  \node[above=0.20cm,font=\footnotesize] at (6.55,4.62)
    {total rank $R = K\,r$, in $K$ contiguous blocks};

  % B: the columns that pair with each block, drawn the same height as A, with a visible
  % gap and the product dot inside it.
  \foreach \i in {1,2,3,4} {
    \draw[line width=0.8pt,fill=white] ({9.30+0.40*(\i-1)},3.90) rectangle ++(0.40,0.72);
  }
  \node[below=0.04cm,font=\small] at (10.10,3.90) {$B$};
  \node[font=\small] at (9.02,4.26) {$\cdot$};

  % ---- band 3 (y = 0.80 to 2.30): the three operations, one box each, with gaps -------
  \node[draw,line width=0.7pt,minimum width=4.4cm,minimum height=1.50cm,fill=white,
        align=left,inner sep=5pt,font=\footnotesize] (write) at (2.40,1.55)
    {\textbf{\small write}\\ gradient steps touch\\ block $k$ of $A$ only};
  \node[draw,line width=0.7pt,minimum width=4.4cm,minimum height=1.50cm,fill=white,
        align=left,inner sep=5pt,font=\footnotesize] (read) at (7.50,1.55)
    {\textbf{\small read}\\ a mask selects block $k$,\\ so nothing enters the context};
  \node[draw,line width=0.7pt,minimum width=4.4cm,minimum height=1.50cm,fill=white,
        align=left,inner sep=5pt,font=\footnotesize] (erase) at (12.60,1.55)
    {\textbf{\small erase}\\ restore $A_k$ and zero $B_k$:\\ bit-identical to never written};

  % One connector from the shaded block, forking to the three boxes. The fork sits in the
  % empty band between the strip's labels (y=3.85) and the boxes' tops (y=2.30).
  \draw[line width=0.7pt,dash pattern=on 2.2pt off 1.6pt] (5.25,3.86) -- (5.25,2.95);
  \draw[line width=0.7pt,dash pattern=on 2.2pt off 1.6pt] (2.40,2.95) -- (12.60,2.95);
  \foreach \x in {2.40,7.50,12.60} {
    \draw[-{Latex[length=2.0mm]},line width=0.7pt,dash pattern=on 2.2pt off 1.6pt]
      (\x,2.95) -- (\x,2.32);
  }

\end{tikzpicture}
\caption{The side path. One low-rank adapter per projection layer, whose total rank $R$ is
split into $K$ contiguous blocks; block $k$ \emph{is} slot $k$. Writing touches only that
block, reading selects it with a mask so the content never enters the context, and erasing
restores the block's initial factor and zeroes its counterpart, which makes forgetting
verifiable rather than approximate. The distinction from work that organises a
\emph{collection} of adapters into slots is that here the slots are rank blocks inside a
single adapter, and that is what makes the three operations independently checkable.}
\label{fig:mechanism}
\end{figure}
